\chapter{2022年天津卷}

\section{单选题}

\begin{problem}
设全集 \(U=\{-2,-1,0,1,2\}\),集合 \(A=\{0,1,2\}\),\(B=\{-1,2\}\),则 \(A\cap\left(C_UB\right)=\)(\quad)

\choices
    {\(\{0,1\}\)}
    {\(\{0,1,2\}\)}
    {\(\{-1,1,2\}\)}
    {\(\{-1,0,1,2\}\)}
\end{problem}

\begin{problem}
``\(x\) 为整数''是``\(2x+1\) 为整数''的(\quad)

\choices
    {充分不必要条件}
    {必要不充分条件}
    {充要条件}
    {既不充分也不必要条件}
\end{problem}

\begin{problem}
函数 \(y=\dfrac{|x^2-1|}{x}\) 的图象大致为(\quad)

\choices
    {\choicebitmap[width=0.15\paperwidth]{img/2022/tianjin/q3_option_a.png}}
    {\choicebitmap[width=0.15\paperwidth]{img/2022/tianjin/q3_option_b.png}}
    {\choicebitmap[width=0.15\paperwidth]{img/2022/tianjin/q3_option_c.png}}
    {\choicebitmap[width=0.15\paperwidth]{img/2022/tianjin/q3_option_d.png}}
\end{problem}

\begin{problem}
将 1916 年到 2015 年的全球年平均气温(单位:\({}^{\circ}\mathrm C\)),共 100 个数据,分成 6 组:\([13.55,13.75)\),\([13.75,13.95)\),\([13.95,14.15)\),\([14.15,14.35)\),\([14.35,14.55)\),\([14.55,14.75]\),并整理得到如下的频率分布直方图,则全球年平均气温在区间 \([14.35,14.75]\) 内的有(\quad)

\bitmapfigure[max width=7.2cm]{img/2022/tianjin/q4_fig1.png}

\choices
    {\(22\text{ 年}\)}
    {\(23\text{ 年}\)}
    {\(25\text{ 年}\)}
    {\(35\text{ 年}\)}
\end{problem}

\begin{problem}
设 \(a=2^{0.7}\),\(b=\left(\dfrac{1}{3}\right)^{0.7}\),\(c=\log_2\dfrac{1}{3}\),则 \(a\),\(b\),\(c\) 的大小关系为(\quad)

\choices
    {\(a<b<c\)}
    {\(c<a<b\)}
    {\(b<c<a\)}
    {\(c<b<a\)}
\end{problem}

\begin{problem}
化简 \(\left(2\log_4 3+\log_8 3\right)\left(\log_3 2+\log_9 2\right)=\)(\quad)

\choices
    {1}
    {\(\dfrac{5}{4}\)}
    {2}
    {\(\dfrac{5}{2}\)}
\end{problem}

\begin{problem}
已知双曲线 \(\dfrac{x^2}{a^2}-\dfrac{y^2}{b^2}=1\)(\(a>0\),\(b>1\))的左,右焦点分别为 \(F_1\),\(F_2\),抛物线 \(y^2=4\sqrt{5}x\) 的准线经过 \(F_1\),且与双曲线的一条渐近线交于点 \(A\). 若 \(\angle F_1F_2A=\dfrac{\pi}{4}\),则双曲线的方程为(\quad)

\choices
    {\(\dfrac{x^2}{16}-\dfrac{y^2}{4}=1\)}
    {\(\dfrac{x^2}{4}-\dfrac{y^2}{16}=1\)}
    {\(\dfrac{x^2}{4}-y^2=1\)}
    {\(x^2-\dfrac{y^2}{4}=1\)}
\end{problem}

\begin{problem}
``十字歇山顶''是中国古代建筑屋顶的经典样式之一,图 1 中的故宫角楼的顶部即为十字歇山顶. 其上部可视为由两个相同的直三棱柱交叠而成的几何体(图 2). 这两个直三棱柱有一个公共侧面 \(ABCD\),在底面 \(BCE\) 中,若 \(BE=CE=3\),\(\angle BEC=120^{\circ}\),则该几何体的体积为(\quad)

\examfiguregroup[float=inline]{img/2022/tianjin/q8_fig1.png}{%
    \begin{tabular}{cc}
    \bitmapinclude[max width=6.0cm]{img/2022/tianjin/q8_fig1.png} &
    \bitmapinclude[max width=6.0cm]{img/2022/tianjin/q8_fig2.png}
    \end{tabular}%
}

\choices
    {\(\dfrac{27}{2}\)}
    {\(\dfrac{27\sqrt{3}}{2}\)}
    {27}
    {\(27\sqrt{3}\)}
\end{problem}

\begin{problem}
关于函数 \(f(x)=\dfrac{1}{2}\sin 2x\),给出下列结论:

\begin{circlelist}
\item \(f(x)\) 的最小正周期是 \(2\pi\);
\item \(f(x)\) 在区间 \(\left[-\dfrac{\pi}{4},\dfrac{\pi}{4}\right]\) 上单调递增;
\item 当 \(x\in\left[-\dfrac{\pi}{6},\dfrac{\pi}{3}\right]\) 时,\(f(x)\) 的取值范围为 \(\left[-\dfrac{\sqrt{3}}{4},\dfrac{\sqrt{3}}{4}\right]\);
\item \(f(x)\) 的图象可以由函数 \(g(x)=\dfrac{1}{2}\sin\left(2x+\dfrac{\pi}{4}\right)\) 的图象向左平移 \(\dfrac{\pi}{8}\) 个单位长度得到.
\end{circlelist}
其中正确结论的个数为(\quad)

\choices
    {1}
    {2}
    {3}
    {4}
\end{problem}

\section{填空题}

\begin{problem}
已知 \(\i\) 是虚数单位,化简 \(\dfrac{11-3\i}{1+2\i}\) 的结果为\fillinblank{}.
\end{problem}

\begin{problem}
在 \(\left(\sqrt{x}+\dfrac{3}{x^2}\right)^5\) 的展开式中,常数项是\fillinblank{}.
\end{problem}

\begin{problem}
若直线 \(x-y+m=0\)(\(m>0\))被圆 \((x-1)^2+(y-1)^2=3\) 截得的弦长等于 \(m\),则 \(m\) 的值为\fillinblank{}.
\end{problem}

\begin{problem}
现有 52 张扑克牌(去掉大小王),每次取一张,取后不放回,则两次都抽到 \(A\) 的概率为\fillinblank{};在第一次抽到 \(A\) 的条件下,第二次也抽到 \(A\) 的概率为\fillinblank{}.
\end{problem}

\begin{problem}
在 \(\triangle ABC\) 中,点 \(D\) 为 \(AC\) 的中点,点 \(E\) 满足 \(\overrightarrow{CB}=2\overrightarrow{BE}\). 记 \(\overrightarrow{CA}=\bs{a}\),\(\overrightarrow{CB}=\bs{b}\),用 \(\bs{a}\),\(\bs{b}\) 表示 \(\overrightarrow{DE}=\)\fillinblank{};若 \(\overrightarrow{AB}\perp\overrightarrow{DE}\),则 \(\angle ACB\) 的最大值为\fillinblank{}.
\end{problem}

\begin{problem}
设 \(a\in\R\),对任意实数 \(x\),记 \(f(x)\) 表示 \(|x|-2\),\(x^2-ax+3a-5\) 中的较小者. 若函数 \(f(x)\) 至少有 3 个零点,则 \(a\) 的取值范围为\fillinblank{}.
\end{problem}

\section{解答题}

\begin{problem}
在 \(\triangle ABC\) 中,角 \(A\),\(B\),\(C\) 所对的边分别为 \(a\),\(b\),\(c\). 已知 \(a=\sqrt{6}\),\(b=2c\),\(\cos A=-\dfrac{1}{4}\).

\begin{enumerate}
\item 求 \(c\) 的值;
\item 求 \(\sin B\) 的值;
\item 求 \(\sin(2A-B)\) 的值.
\end{enumerate}
\end{problem}

\begin{problem}
如图,在直三棱柱 \(ABC-A_1B_1C_1\) 中,\(AC\perp AB\),点 \(D\),\(E\),\(F\) 分别为 \(A_1B_1\),\(AA_1\),\(CD\) 的中点,\(AB=AC=AA_1=2\).

\begin{enumerate}
\item 求证:\(EF\parallel\) 平面 \(ABC\);
\item 求直线 \(BE\) 与平面 \(CC_1D\) 所成角的正弦值;
\item 求平面 \(A_1CD\) 与平面 \(CC_1D\) 所成二面角的余弦值.
\end{enumerate}

\bitmapfigure[max width=4.8cm]{img/2022/tianjin/q17_fig1.png}
\end{problem}

\begin{problem}
设 \(\{a_n\}\) 为等差数列,\(\{b_n\}\) 为等比数列,且 \(a_1=b_1=a_2-b_2=a_3-b_3=1\).

\begin{enumerate}
\item 求 \(\{a_n\}\) 和 \(\{b_n\}\) 的通项公式;
\item 记 \(\{a_n\}\) 的前 \(n\) 项和为 \(S_n\),求证:\(\left(S_{n+1}+a_{n+1}\right)b_n=S_{n+1}b_{n+1}-S_nb_n\);
\item 求 \(\displaystyle\sum_{k=1}^{2n}\left[a_{k+1}-(-1)^ka_k\right]b_k\).
\end{enumerate}
\end{problem}

\begin{problem}
椭圆 \(\dfrac{x^2}{a^2}+\dfrac{y^2}{b^2}=1\)(\(a>b>0\))的右焦点为 \(F\),右顶点为 \(A\),上顶点为 \(B\),且满足 \(\dfrac{|BF|}{|AB|}=\dfrac{\sqrt{3}}{2}\).

\begin{enumerate}
\item 求椭圆的离心率 \(e\);
\item 直线 \(l\) 与椭圆有唯一公共点 \(M\),与 \(y\) 轴相交于 \(N\)(\(N\) 异于 \(M\)). 记 \(O\) 为坐标原点,若 \(|OM|=|ON|\),且 \(\triangle OMN\) 的面积为 \(\sqrt{3}\),求椭圆的标准方程.
\end{enumerate}
\end{problem}

\begin{problem}
已知 \(a\),\(b\in\R\),函数 \(f(x)=\e^x-a\sin x\),\(g(x)=b\sqrt{x}\).

\begin{enumerate}
\item 求函数 \(y=f(x)\) 在点 \((0,f(0))\) 处的切线方程;
\item 若 \(y=f(x)\) 和 \(y=g(x)\) 有公共点,
\begin{enumerate}
\item 当 \(a=0\) 时,求 \(b\) 的取值范围;
\item 求证:\(a^2+b^2>\e\).
\end{enumerate}
\end{enumerate}
\end{problem}

